\documentclass[10pt,a4paper]{article}

\usepackage[utf8]{inputenc}
\usepackage[T1]{fontenc}
\usepackage[french]{babel}
\usepackage{amsmath,amssymb}
\usepackage[top=1cm,bottom=1.2cm,left=1.5cm,right=1.5cm]{geometry}
\usepackage{hyperref}
\hypersetup{colorlinks=true,linkcolor=black,citecolor=blue,urlcolor=blue}

\title{Analyse dimensionnelle et dynamique d'accumulation}
\author{Salomé Fonvielle}
\date{}

\begin{document}
\maketitle

\noindent\emph{Question traitée.} Quelle est l'unité physique de $\tau_{\mathrm{ARF}}$
et celle du seuil $\lambda$ ? En quoi leur comparaison directe contourne-t-elle la
nature séquentielle du détecteur ?

\paragraph{Les deux unités.} Dans un flux, il n'existe pas d'horloge physique : la
seule variable temporelle disponible est l'indice de l'exemple qui arrive, et
$\tau_{\mathrm{ARF}}$ se mesure donc en \emph{observations}. Le seuil, lui, se lit sur
la récurrence du détecteur cumulatif\footnote{Équation~(3) du manuscrit, section~III-D
(« The Starvation Effect »), notation reprise telle quelle. Le protocole R2 crée le
détecteur à la rupture, d'où $S_{\tau^*} = 0$ ; la Proposition~3 du manuscrit conserve
au contraire $S_{\tau^*} = s_0 \ge 0$, poser $s_0 = 0$ est ici une spécialisation du
protocole.}
\begin{equation}
  S_t = \max\bigl(0,\, S_{t-1} + e_t - p_0 - \delta_P\bigr)\,,
  \label{eq:cusum}
\end{equation}
où $e_t \in \{0,1\}$ indique une erreur au pas $t$, $p_0$ le taux d'erreur normal avant
la rupture et $\delta_P$ une tolérance retranchée à chaque pas. Comme $e_t$ vaut $0$ ou
$1$, une erreur fait monter $S$ de près d'une unité tandis qu'un succès le fait à peine
redescendre : à la tolérance près, $S_t$ compte les erreurs accumulées au-delà du socle.
Puisque $\lambda$ n'apparaît que dans la comparaison $S_t \ge \lambda$, il se lit
nécessairement dans la même unité --- $\lambda = 50$ ne signifie pas « attendre 50 pas »,
mais « attendre qu'environ 50 erreurs excédentaires se soient accumulées ». Les deux
quantités sont donc des comptes de natures différentes : écrire
$\tau_{\mathrm{ARF}} < \lambda$ n'a aucun sens dimensionnel.

\paragraph{Ce que la comparaison directe escamote.} Comparer les deux revient à traiter
$\lambda$ comme une durée : « il faut $\lambda$ erreurs, donc tant de pas ». Ce raccourci
suppose que la preuve s'accumule à un rythme fixe --- sans un nombre d'erreurs par pas,
on ne passe pas d'une unité à l'autre --- et l'hypothèse est contenue dans le geste même
de comparer, qu'on écrive la conversion ou non. Or ce rythme n'existe pas : le détecteur
reçoit les erreurs une par une et ne produit pas un nombre, mais une date, celle où la
trajectoire de $S_t$ franchit $\lambda$. Trois traits de cette trajectoire disparaissent
lorsqu'on la remplace par une montée à pente constante.

\emph{Le rythme ralentit.} Juste après la rupture, l'incrément du compteur a pour
espérance $\Delta e - \delta_P$ ; mais la forêt se répare, l'excès d'erreur redescend
vers le socle et ce gain s'épuise. Le manuscrit suppose au contraire un excès constant à
$\Delta e$ pendant tout le transitoire, puis nul --- un rectangle, dont la hauteur est
justement le rythme fixe qu'exige la comparaison.

\emph{La preuve acquise peut être détruite.} Quand $S_t$ passerait sous zéro, le
$\max(0,\cdot)$ de~\eqref{eq:cusum} le ramène à $0$ : ce n'est pas une mise en pause,
tout ce qui avait été accumulé est perdu. Les mêmes erreurs, groupées juste après la
rupture ou dispersées sur l'horizon, ne donnent donc pas la même date d'alarme, alors
qu'un rythme moyen ne fait aucune différence entre les deux. C'est le sens précis de
« séquentiel » : la date dépend de l'ordre d'arrivée, pas seulement du total.

\emph{L'alarme peut ne pas venir.} Sur l'horizon $H = 2\,000$ de la campagne, le seuil
peut n'être jamais franchi\footnote{Au-delà de l'horizon le franchissement finit par
survenir, mais son délai moyen croît exponentiellement avec $\lambda$ (manuscrit,
Remarque~4) : l'alarme ne porte alors plus sur la dérive.}.

\paragraph{Conclusion.} Un compte d'erreurs et un compte d'observations ne se comparent
pas en postulant un rythme. L'homogénéité se rétablit dans l'autre sens : en ramenant les
deux grandeurs dans l'unité de preuve, où la comparaison se fait seuil contre évidence
accumulée, sur chaque exécution et sans moyenne. C'est l'objet de la statistique
culminante $S_{\max}(H)$ et du certificat déterministe qui en découle.

% NOTE(coordination) : la conversion tau*_det = lambda/(Delta e - delta_P) et la
% substitution du temps d'arret par cette constante relevent de la question A
% (Proposition 9, versant theorique) : volontairement non traitees ici.

\end{document}
